\chapter*{Nota metodológica: reprodutibilidade}
\addcontentsline{toc}{chapter}{Nota metodológica: reprodutibilidade}
\label{chap:nota_metodologica}

Esta nota reúne as escolhas de implementação que permitem reproduzir os
resultados. As definições das variáveis estão no
Apêndice~\ref{app:variaveis}; os scripts, no Apêndice~\ref{app:scripts}; e a
ordem de execução, no Apêndice~\ref{app:pastas}. As limitações da análise
estão na Seção~\ref{sec:concl-limitacoes}.

\section*{Dados e amostras}

Os preços de fechamento diários do Bitcoin (par BTCUSDT) vêm da API pública da
Binance, e o índice DVOL, da API pública da Deribit. A série de preços não
tinha os dias de 1º a 31 de março de 2023, que a dissertação completa com as
velas diárias da mesma fonte; o pipeline interrompe a execução se faltar
qualquer data. A amostra completa tem 1.806 observações, de 24/03/2021 a
03/03/2026, e serve à análise descritiva e ao teste de H1; a amostra de
modelagem tem 1.776, de 23/04/2021 a 03/03/2026, e serve aos
Capítulos~\ref{chap:previsao_bvrp} a~\ref{chap:retornos_nao_linear}
(Seção~\ref{sec:met-amostras}).

\section*{Informação disponível em \texorpdfstring{$t$}{t}}

As variáveis explicativas de cada data são conhecidas no fim do próprio dia.
O BVRP de $t$ só é conhecido em $t+30$ e nunca entra como variável
explicativa. A dissertação reestima os modelos a cada 30 dias, em 33 origens,
e o treino de cada origem $t$ usa apenas as datas $s \le t - h$; a validação
cruzada, com cinco dobras expansivas, aplica o mesmo embargo dentro do treino.
O corte principal entre os regimes usa só a primeira janela de estimação. Um
teste placebo, com o alvo embaralhado dentro do treino, confere que o
procedimento não usa informação do período de teste
(Capítulo~\ref{chap:previsao_bvrp}).

\section*{Modelos e inferência}

O MQO, o Ridge, o LASSO e a floresta aleatória usam a biblioteca
\texttt{scikit-learn}, e o XGBoost, a biblioteca \texttt{xgboost}; os modelos
lineares usam as variáveis padronizadas com a média e o desvio-padrão do
treino de cada origem. As regressões usam o MQO do \texttt{statsmodels} com
erro-padrão de Newey--West: núcleo de Bartlett com $h+1$ defasagens, sem
correção de amostra pequena e com valor-$p$ pela distribuição normal. O
bootstrap em blocos em dois níveis usa 999 reamostragens e blocos de 60 dias,
com blocos de 30 e de 90 dias como robustez (Seção~\ref{sec:met-inferencia}).
O \textit{model confidence set} do Capítulo~\ref{chap:previsao_bvrp} usa 9.999
reamostragens em blocos de 60 dias, com 30 e 90 dias como sensibilidade
(Seção~\ref{sec:met-avaliacao}).

\section*{Reprodutibilidade}

As sementes são fixas: 42 nos modelos em árvore e nas permutações do placebo,
e 20261001 nas reamostragens do bootstrap e do \textit{model confidence set}. As
tabelas e as figuras do texto
são gravadas pelos scripts, sem edição manual; os geradores recalculam as
estatísticas reportadas a partir das saídas gravadas e interrompem a execução
se a diferença passar de $10^{-9}$. O ambiente usado tem Python 3.14.4,
\texttt{pandas} 3.0.2, \texttt{NumPy} 2.4.4, \texttt{statsmodels} 0.14.6,
\texttt{scikit-learn} 1.8.0 e \texttt{xgboost} 3.2.0; o arquivo
\texttt{requirements.txt} do repositório fixa essas versões e as das demais
bibliotecas.

\section*{Disponibilidade}

O código e os dados estão disponíveis mediante solicitação ao autor.
